\section{Mencatat lahan}

Buka menu \textbf{Manajemen Lahan}, lalu pilih \textbf{Tambah lahan} (Gambar~\ref{fig:daftar-lahan}). Isi informasi berdasarkan dokumen resmi dan kondisi lapangan (Gambar~\ref{fig:tambah-lahan}). Isian bertanda bintang (\texttt{*}) wajib dilengkapi.

\guideimage{07-list-management-lahan.png}{Daftar Manajemen Lahan. Gunakan pencarian atau filter untuk menemukan lahan; pilih Tambah lahan untuk mencatat lahan baru.}{fig:daftar-lahan}

\guideimage{01-form-tambah-lahan.png}{Form tambah lahan pada menu Manajemen Lahan. Lengkapi informasi dan lokasi lahan secara berurutan, unggah foto dokumentasi bila ada, lalu pilih Simpan lahan.}{fig:tambah-lahan}

\subsection{Informasi lahan}

\begin{longtable}{|>{\bfseries\raggedright\arraybackslash}p{\dimexpr 0.27\linewidth - 2\tabcolsep - \arrayrulewidth\relax}|p{\dimexpr 0.73\linewidth - 2\tabcolsep - 2\arrayrulewidth\relax}|}
  \tableheaderrow{Isian}{Cara mengisi}
  \endhead
  Nama lahan & Gunakan nama yang singkat dan mudah dikenali, misalnya \textquotedblleft Lahan Cileuksa\textquotedblright\ atau \textquotedblleft Lahan Cariu\textquotedblright. Hindari nama yang sama dalam satu ruang kerja. \\ \hline
  Luas lahan (ha) & Masukkan luas resmi/terverifikasi dalam hektare (maksimal 4 angka desimal). Perhatikan satuan; $1\text{ ha} = 10.000\text{ m}^2$. \\ \hline
  Target volume panen (ton) & Masukkan target total sesuai rencana usaha ($1\text{ ton} = 1.000\text{ kg}$). Jika target dihitung dari $\text{luas} \times \text{produktivitas per hektare}$, pastikan komoditas, satuan, periode, dan angka produktivitas disepakati terlebih dahulu. Form meminta nilai target secara langsung. \\ \hline
  Status penguasaan lahan & Pilih \textbf{Milik sendiri}, \textbf{Sewa}, atau \textbf{Kerja sama} sesuai dokumen. Jika memilih \textbf{Sewa} atau \textbf{Kerja sama}, lengkapi kolom wajib \textbf{Nama pemilik lahan} atau \textbf{Nama mitra kerja sama}. \\ \hline
  PIC lahan & Opsional. Pilih akun Petugas yang telah tersedia. Penugasan dapat dilakukan atau diperbarui kemudian. \\ \hline
\end{longtable}

\subsection{Lokasi dan dokumentasi}

\begin{enumerate}
  \item Tulis alamat yang membantu menemukan lokasi, termasuk desa/dusun atau petunjuk akses. Hindari memasukkan nama pribadi, nomor telepon, atau informasi rahasia karena alamat dikirim ke layanan pencarian peta.
  \item Pilih \textbf{Provinsi}, kemudian \textbf{Kota/kabupaten}, lalu \textbf{Kecamatan}. Pilihan wilayah berikutnya aktif setelah pilihan sebelumnya tersedia.
  \item Pada bagian \textbf{Titik lokasi lahan (opsional)}, sistem dapat menampilkan pilihan titik lokasi berdasarkan alamat dan wilayah yang diisi. Anda juga dapat mengeklik peta untuk menaruh pin, lalu menggeser penanda jika posisinya perlu dikoreksi. Pin menunjukkan satu titik referensi, bukan batas luas lahan.
  \item Pada bagian \textbf{Dokumentasi lahan}, Anda dapat mengunggah maksimal 1 \textbf{Foto lahan (opsional)} berformat JPEG, PNG, atau WebP dengan batas ukuran 15~MB (ukuran foto disesuaikan otomatis sebelum diunggah).
  \item Pilih \textbf{Simpan lahan}, lalu buka detailnya untuk memastikan luas, wilayah, status, PIC, titik lokasi, dan foto sudah benar.
\end{enumerate}
